\documentclass[oneside,a4paper,14pt]{extarticle}
\usepackage[a4paper,top=20mm,bottom=20mm,left=30mm,right=10mm]{geometry}
\usepackage{amsmath, amssymb, amsfonts, amsthm}
\usepackage[T2A]{fontenc}
\usepackage[utf8]{inputenc}
\usepackage[russian]{babel}
\usepackage{textcomp}
\usepackage{indentfirst}
\usepackage{graphicx}
\usepackage{float}
\usepackage{wrapfig}
\usepackage{caption}
\usepackage{tocloft}
\renewcommand{\cftsecleader}{\cftdotfill{\cftdotsep}}
\usepackage[breaklinks]{hyperref}
\usepackage{titlesec}
\usepackage{fancyvrb}
\usepackage{xcolor}
\titleformat{\section}{\normalsize\bfseries}{\thesection}{1em}{}
\titleformat{\subsection}{\normalsize\bfseries}{\thesubsection}{1em}{}
\titleformat{\subsubsection}{\normalsize\bfseries}{\thesubsubsection}{1em}{}
\renewcommand\baselinestretch{1.33}\normalsize
\setlength{\parindent}{1.25cm}
z
\begin{document}
\newpage\thispagestyle{empty}
\newpage
\setcounter{page}{1}
\tableofcontents
\newpage

\section*{Задание}
\addcontentsline{toc}{section}{Задание}
Задание для шестого варианта:\\
«Сдвигатель» 8-разрядного двоичного числа, содержащего в старшем (седьмом) разряде знак числа, на N $\le$ 7 разрядов в сторону младших разрядов (сдвиг арифметический)

\section*{Анализ задания по разработке функциональной модели}
\addcontentsline{toc}{section}{Анализ задания по разработке функциональной модели}

Разработать функциональную модель двоичного сдвигателя.\\

На первый вход подаётся 8-разрядный двоичный код\\$X=x_0,x_1,x_2,x_3,x_4,x_5,x_6,x_7$,
	а на второй вход 3-разрядный двоичный код переноса $E=e_0,e_1,e_2$.\\

	При $e_0,e_1,e_2=0$ на информационный выход $Y=y_0,y_1,y_2,y_3,y_4,y_5,y_6,y_7$ код поступает без изменений.\\

	При $x_0=1,x_1=0,x_2=1,x_3=0,x_4=0,x_5=1,x_6=0,x_7=0$ (10100100) $e_2=1$ на информационном выходе формируется сдвинутый вправо на 1 бит 8-разрядный двоичный код:\\
$y_0=1,y_1=1,y_2=0,y_3=1,y_4=0,y_5=0,y_6=1,y_7=0$\\
	(11010010)

	При этом выход $y_7$ всегда будет равен $x_7$ и будет определять бит, которым будут заменяться младшие разряды.

\section*{Вывод}
\addcontentsline{toc}{section}{Вывод}
В ходе выполнения лабораторной работы была разработана схема булевых функций, описывающая работу сдвигателя. На основе схемы булевых функций были составлены их представления и составлена экранная форма узла для экспериментального исследования инструментами табличного процессора MS Excel. В ходе экспериментального исследования составлена таблица результатов, которая отражает работу узла без ошибок.

\end{document}